% ======================================================================
% Section 5 -- Results I: Explaining predictions (~2.4 pp)
% Bold-first-sentence paragraphs (\expclaim). Numbers live in tables.
% No em dashes. Multi-seed error bars everywhere.
% ======================================================================
\section{Results I: Explaining Predictions}
\label{sec:predictions}

\textbf{Setup.} We use the four GraphXAI molecular benchmarks with
annotated ground-truth motifs \citep{agarwal2023graphxai}, the
TopoBench datasets NCI1 and PROTEINS, which carry no explanation
annotations, and MANTRA \citep{ballester2025mantra}, a native
simplicial dataset. Every comparison uses identical training splits,
the full ground-truth-positive test population, and three training
seeds; every entry reports mean $\pm$ standard deviation over seeds.
These experiments realize the framework on
cellular complexes (the molecular benchmarks, lifted) and simplicial
complexes (MANTRA). Explanation accuracy is measured by GEA: select $b$ cells, where $b$ is
the size of a ground-truth candidate, and report the Jaccard overlap
with that candidate, maximized over candidates. Faithfulness is
measured by deletion-AUC: detach cells in score order and integrate the
predicted-class probability, so lower is better. Full protocols are in
\cref{app:experimental}.

\subsection{Validation at Rank 1: Graph Benchmarks}
\label{sec:predictions-rank1}

On graphs the participation game reduces to induced-subgraph removal
(\cref{sec:framework}). This lets us compare \method against GNN
explainers on their own terms. All methods explain the same trained,
parameter-matched GCN and GIN subjects.

% AUTO-GENERATED by scripts/make_paper_tables.py -- do not edit
\begin{table}[t]
\centering
\small
\caption{\textbf{Graph benchmarks (rank-1 validation).} Node-level explanation accuracy (GEA, mean $\pm$ sd over 3 training seeds) on the full ground-truth-positive test splits. All methods explain the same parameter-matched GNN subjects on identical splits; per method, the better of the GCN/GIN subjects is shown (both in \cref{app:tables}). \method{} explanations cost $\sim$0.04\,s per instance vs.\ $\sim$20\,s for SubgraphX on identical hardware.}
\label{tab:graph-benchmarks}
\begin{tabular}{lcccc}
\toprule
 & Benzene & FluorideCarbonyl & Mutagenicity & AlkaneCarbonyl \\
\midrule
\method{} (attribution) & 0.859 $\pm$ 0.009 & \cellcolor{black!10}0.507 $\pm$ 0.047 & \cellcolor{black!10}\textbf{0.798 $\pm$ 0.124} & \cellcolor{black!10}\textbf{0.500 $\pm$ 0.081} \\
\method{} (explanation) & \cellcolor{black!10}\textbf{0.969 $\pm$ 0.009} & 0.340 $\pm$ 0.011 & 0.480 $\pm$ 0.135 & 0.354 $\pm$ 0.125 \\
SubgraphX & 0.796 $\pm$ 0.029 & 0.439 $\pm$ 0.107 & 0.344 $\pm$ 0.036 & 0.296 $\pm$ 0.010 \\
PGExplainer & 0.642 $\pm$ 0.044 & \cellcolor{black!10}\textbf{0.555 $\pm$ 0.067} & 0.290 $\pm$ 0.102 & 0.205 $\pm$ 0.160 \\
GNNExplainer & 0.443 $\pm$ 0.054 & 0.210 $\pm$ 0.014 & 0.158 $\pm$ 0.023 & 0.162 $\pm$ 0.049 \\
Random & 0.250 $\pm$ 0.000 & 0.219 $\pm$ 0.000 & 0.132 $\pm$ 0.000 & 0.173 $\pm$ 0.000 \\
\bottomrule
\end{tabular}
\end{table}


\expclaim{\method matches or outperforms every baseline on the
GraphXAI benchmarks.} \Cref{tab:graph-benchmarks} reports both
outputs. The explanation output leads on Benzene and the attribution
output leads on Mutagenicity and Alkane-Carbonyl. PGExplainer retains
the best mean on Fluoride-Carbonyl, within overlapping error bars of
our attribution output. The two outputs are complementary: attribution
measures how much the prediction depends on each cell, and wins where
the motif is necessary; explanation measures joint sufficiency, and
wins where the motif alone preserves the prediction.

\expclaim{The explanation output is orders of magnitude cheaper than
search-based explanation.} Greedy pruning evaluates each candidate
removal once per step, one forward pass per coalition. SubgraphX pays
roughly two hundred forward passes per coalition it scores, because its
Monte-Carlo reward samples coalition contexts. Measured wall-clock
costs on identical hardware appear in the caption of
\cref{tab:graph-benchmarks}.

\expclaim{Both removal semantics and search direction matter, and both
choices are ablated.} Replacing detachment with feature zero-filling,
the convention of SubgraphX, reproduces SubgraphX's behavior on GCN
subjects and degrades substantially on GIN subjects
(\cref{tab:by-subject}). Growing coalitions from single cells instead of
pruning from the full complex fails: early evaluations score
coalitions far outside the training distribution
(\cref{app:experimental}).

\expclaim{Where a baseline wins, the game explains why.} Define the
alignment gap $\Delta = v(S^{*}_{b}) - v(S_{\mathrm{gt}})$, where
$S_{\mathrm{gt}}$ is the ground-truth motif as a coalition. On
Fluoride-Carbonyl the explained GIN scores non-motif coalitions above
the annotated motif on 96\% of candidates, so every faithful explainer
is capped there. GEA evaluates the composition of explainer and model
against human annotations. $\Delta$ separates the two factors at the
cost of two extra game evaluations per instance, and no baseline can
compute it.

\subsection{Higher-Order Domains}
\label{sec:predictions-higher}

We now explain TNNs on lifted complexes, where faces are rank-2
players. The explained model in every experiment is the
full-vocabulary GCCN, which uses all eleven candidate neighborhoods
and involves no selection step.

\expclaim{Standard feature liftings cache structure into features, and
models trained on them are harder to explain.} The standard lifting
computes a face's features by summing its nodes' features at
preprocessing time. The face then carries information about its
boundary even after the boundary is detached, so the counterfactuals
of any removal-based explainer leak. Trained models exploit the cache:
their sufficient coalitions keep the face while dropping its
boundary (\cref{app:experimental}), and their alignment gap is
positive on nearly all candidates.

% AUTO-GENERATED by scripts/make_paper_tables.py -- do not edit
\begin{table}[t]
\centering
\small
\caption{\textbf{Lifted topological benchmarks.} Node-projected GEA (mean $\pm$ sd, 3 seeds, full populations) for both \method{} outputs on the same full-vocabulary GCCN trained under two feature liftings. The standard lifting sums boundary features into higher-rank cells; the structural lifting uses constant features. Bold marks the best entry per column. Subject accuracies: Benzene 0.942, FluorideCarbonyl 0.947, Mutagenicity 0.976, AlkaneCarbonyl 0.967; alignment-ceiling share $\Pr[v(S^{*})>v(S_{\mathrm{gt}})]$: Benzene 0.74, FluorideCarbonyl 0.93, Mutagenicity 0.97, AlkaneCarbonyl 0.99.}
\label{tab:lifted}
\begin{tabular}{lcccc}
\toprule
 & Benzene & FluorideCarbonyl & Mutagenicity & AlkaneCarbonyl \\
\midrule
attribution, structural lifting & 0.471 $\pm$ 0.018 & \cellcolor{black!10}0.491 $\pm$ 0.035 & 0.326 $\pm$ 0.065 & \cellcolor{black!10}0.333 $\pm$ 0.097 \\
attribution, standard lifting & 0.459 $\pm$ 0.192 & \cellcolor{black!10}\textbf{0.516 $\pm$ 0.089} & 0.257 $\pm$ 0.077 & \cellcolor{black!10}0.348 $\pm$ 0.043 \\
explanation, structural lifting & \cellcolor{black!10}\textbf{0.725 $\pm$ 0.105} & 0.426 $\pm$ 0.027 & \cellcolor{black!10}\textbf{0.436 $\pm$ 0.081} & \cellcolor{black!10}0.362 $\pm$ 0.130 \\
explanation, standard lifting & 0.509 $\pm$ 0.188 & 0.372 $\pm$ 0.004 & 0.289 $\pm$ 0.117 & \cellcolor{black!10}\textbf{0.364 $\pm$ 0.053} \\
\bottomrule
\end{tabular}
\end{table}


\expclaim{A structure-only lifting improves accuracy and explanation
quality together.} The structural lifting assigns constant features to
all cells above rank 0, the convention MANTRA uses natively, so
structure exists only in the incidences and must be computed by
message passing. \Cref{tab:lifted} compares the same architecture
under both liftings: the structural lifting raises GEA, lowers the
alignment-ceiling share, and trains to higher accuracy. Explanations
of the structural models recover the face's boundary, and the face
itself appears in the explanation on 92\% of Benzene instances.
A cell-level pointer of this kind is not expressible by any node-based
method.

% AUTO-GENERATED by scripts/make_paper_tables.py -- do not edit
\begin{table}[t]
\centering
\small
\caption{\textbf{Faithfulness without ground truth.} Deletion-AUC (lower is better; predicted-class probability as cells are detached in score order; mean $\pm$ sd over 3 training seeds, 50 pre-declared test complexes each) on TopoBench datasets with no explanation annotations. Subjects: full-vocabulary GCCNs over the structural lift; explanations target the model's own prediction.}
\label{tab:nogt}
\begin{tabular}{lcccc}
\toprule
 & NCI1 & NCI109 & MUTAG & PROTEINS \\
\midrule
\method{} & \cellcolor{black!10}\textbf{0.218 $\pm$ 0.005} & \cellcolor{black!10}\textbf{0.203 $\pm$ 0.044} & \cellcolor{black!10}\textbf{0.175 $\pm$ 0.043} & \cellcolor{black!10}\textbf{0.338 $\pm$ 0.067} \\
Occlusion & 0.254 $\pm$ 0.021 & \cellcolor{black!10}0.240 $\pm$ 0.038 & 0.231 $\pm$ 0.006 & 0.431 $\pm$ 0.065 \\
Gradient$\times$Input & 0.552 $\pm$ 0.016 & 0.507 $\pm$ 0.049 & 0.499 $\pm$ 0.073 & 0.638 $\pm$ 0.047 \\
Random & 0.545 $\pm$ 0.004 & 0.491 $\pm$ 0.035 & 0.481 $\pm$ 0.014 & 0.666 $\pm$ 0.058 \\
\bottomrule
\end{tabular}
\end{table}


\expclaim{Without ground truth, \method is the most faithful method
and reports how much topology a model uses.} On NCI1 and PROTEINS
there are no annotated motifs, which is the common case in practice.
Explanations target the model's own prediction
(\cref{sec:framework-game1}). \Cref{tab:nogt} shows deletion-AUC
against occlusion, gradient-times-input, and random baselines on the
same pre-declared test samples. The attribution output also yields a
diagnostic no baseline provides: the share of Shapley mass per rank.
On NCI1 the trained models place roughly half of their credit on
rank-1 cells and less than a tenth on rank-2 cells, a direct readout
of how much higher-order structure each model actually uses.

\expclaim{On native triangulations, rank-level attributions separate
how models decide.} On MANTRA's orientability task, per-rank
attribution profiles distinguish model families that reach similar
accuracy through different mechanisms, and the null-player axiom
exposed one published configuration whose higher-rank cells received
exactly zero credit: a topological model that never used its topology.
\method is also the most faithful method on MANTRA under deletion-AUC.
Full profiles and the multi-complex protocol are in
\cref{app:figures,app:experimental}.
